% Generated draft: requires mathematical and source review.
\documentclass[11pt]{article}
\usepackage[margin=25mm]{geometry}
\usepackage{fontspec}
\setmainfont{Latin Modern Roman}
\usepackage{amsmath,amssymb,mathtools,unicode-math}
\setmathfont{Latin Modern Math}
\usepackage{enumitem,needspace,xurl,hyperref}
\setlength{\parindent}{0pt}
\setlength{\parskip}{5pt}
\setlength{\emergencystretch}{3em}
\newcommand{\N}{\mathbb{N}}
\newcommand{\Z}{\mathbb{Z}}
\newcommand{\Q}{\mathbb{Q}}
\newcommand{\R}{\mathbb{R}}
\newcommand{\C}{\mathbb{C}}
\newcommand{\F}{\mathbb{F}}
\newcommand{\A}{\mathbb{A}}
\newcommand{\PP}{\mathbb{P}}
\newcommand{\E}{\mathbb{E}}
\newcommand{\eps}{\varepsilon}
\newcommand{\dd}{\,\mathrm{d}}
\begin{document}
\section*{UnsolvedMath ID: 30006343}
\textbf{Source problem code:} catalogue entry
\subsection{Definitions and mathematical statement}
Let \(p\) be an odd prime and let \(G\) be a finite group of order \(p^n\), where \(n\geq 0\). For \(g\in G\), define its conjugacy class and breadth by
\[
 g^G=\{x^{-1}gx:x\in G\},\qquad
 p^{\operatorname{br}(g)}=|g^G|.
\]
Thus \(\operatorname{br}(g)\) is a nonnegative integer. Define the breadth of \(G\) by
\[
 \operatorname{br}(G)=\max_{g\in G}\operatorname{br}(g).
\]
For subgroups \(H,K\leq G\), let \([H,K]\) denote the subgroup generated by all commutators \([h,k]=h^{-1}k^{-1}hk\), with \(h\in H\) and \(k\in K\). Define the lower central series by \(\gamma_1(G)=G\) and \(\gamma_{i+1}(G)=[\gamma_i(G),G]\). The nilpotency class \(\operatorname{cl}(G)\) is the least integer \(c\geq 0\) such that \(\gamma_{c+1}(G)=\{1\}\); in particular, \(\operatorname{cl}(\{1\})=0\).

\textbf{Question.} Is it true that, for every odd prime \(p\) and every finite \(p\)-group \(G\),
\[
 \operatorname{cl}(G)\leq\operatorname{br}(G)+1?
\]
The supplied registry labels this question open. Its dated curation reports counterexamples for \(p=2\) and an unresolved odd-prime case; these literature claims have not been independently verified here. The displayed question follows the supplied corrected statement. The supplied original excerpt does not itself display the inequality, so its reconstruction depends on the registry's correction and verification claim.
\subsection{Short English statement}
For every finite group whose order is a power of an odd prime \(p\), must the nilpotency class be at most one more than the largest exponent of \(p\) occurring as a conjugacy-class size? The supplied record labels this question open; this is a statement-only draft requiring human mathematical review.
\subsection{Sources}
\begin{itemize}
\item Local UnsolvedMath problem registry, permanent ID \(30006343\), exact source problem code \texttt{catalogue entry}, title: \emph{Class-Breadth Inequalities for Odd-Prime Groups}. The record gives proposed year 2025 and no named proposer.
\item \emph{Computational Group Theory} (2025), Oberwolfach Reports, pp.\ 1389--1450, DOI \texttt{10.4171/owr/2025/27}. Supplied source report and literature-source label ``OWR report'': \url{https://doi.org/10.4171/owr/2025/27}.
\item Supplied literature source: \emph{Counterexamples to the class-breadth conjecture}, \url{https://doi.org/10.1112/jlms/s2-24.1.113}.
\item Supplied literature source: \emph{On the class-breadth conjecture for p greater than 2}, \url{https://arxiv.org/abs/2606.13423}.
\item The supplied original excerpt attributes a family of counterexample \(2\)-groups to Felsch, Neub\"user, and Plesken in 1987, using reference label [3]. The bibliography entry corresponding to [3] is not supplied. The excerpt also contains the contribution title \emph{Computing the breadths of algebras and T-groups}, attributed to \(\text{\'Oscar Fern\'andez Ayala}\), with joint work credited to Bettina Eick; its precise relationship to the question is not established by the excerpt.
\item Literature assessment dated exactly \texttt{2026-08-22}: the supplied curation states that the inequality has counterexamples for \(p=2\), that the odd-prime case remains open in current primary literature, and that the extracted odd-prime statement is therefore not disproved. This is supplied dated curation, not an independently verified literature assessment.
\item The registry marks the statement \texttt{corrected\_verified} and states that the official report was checked, the question isolated from surrounding exposition, and notation repaired for self-contained publication. These are supplied source-verification claims; no source visit or independent verification is claimed here.
\end{itemize}
\end{document}
